\section{A common tame field and the constant in the type-prime lifting theorem}

% BEGIN Q84T1
\paragraph{Q84T1. Integral descent in the coefficient field.}
Take any of the local data $D$ tied to the usual periods and any test
surjection $g:T_1\twoheadrightarrow T_2$ of stable free lattices for
$G_{\mathbb Q_p}$ in the definition of $\mathrm{admT}$. Its potentially
crystalline representations use a common finite normal extension $K'$,
an inertia type $\Sigma$ and Frobenius as in T84L1. Write
$N\ge2$, $1\le b<H$, $0\in U\subset[-b,H]$, $|U|\le N+1$ for
the jump parameters. Assume $40(HN^2+N^3)<p$. Put $a=(p-1)!$,
$e=p^a-1$, and enlarge the coefficient field $L$ to its compositum
$L'$ with the unramified extension $\mathbb Q_{p^a}$ of degree $a$.
Its residue field contains $\mathbb F_{p^a}$, so its Teichm\"uller units
contain $\mu_e$. Both $a$ and $e$ are prime to $p$.

Write $O,O'$ for the integer rings. If $0\ne d\in O$, $x\in O'$
and $dx\in O$, then $x\in L\cap O'=O$. Thus $O'/O$ is finite
torsion free over the DVR $O$, hence free. There is an $O$-linear
retraction $s:O'\to O$ with $s(1)=1$. Scalar extension preserves
the surjection, its kernel, its periodic characters and its jumps.
For the latter assertion, choose an $L$-basis of $L'$: the period
invariants and the intersections with each filtered period space
are computed coordinate by coordinate. For a finite integral
cocycle $c$ on $T_2$, an identity
$p^d(1\otimes c)=g_*c'+\partial t'$ after extension descends
by $s\otimes1$. The two images are integral, and the extension
of $s$ to $L'\to L$ takes a crystalline period primitive for
$c'$ to one for its image. This proves descent of the lifting
exponent without a factor $[L':L]$.
% END Q84T1

% BEGIN Q84T2
\paragraph{Q84T2. A common tame field for every input lattice.}
Choose $\pi_0^e=-p$ and $\pi_{j+1}^p=\pi_j$, and put
$K_t=\mathbb Q_{p^a}(\pi_0)$. The ordinary inertia group $I$ maps
onto the inertia group of the finite normal extension $K'$.
Indeed a lift of an element acting trivially on the residue field
of $K'$ can be corrected by an element of $G_{K'}$ with the inverse
residue action. Thus characters of the inertia group in the period
datum that are trivial on $G_{K'}$ are distinguished by their
restrictions to $I$. Their restrictions to wild inertia are trivial:
the finite quotient is a $p$-group there and each character has
order prime to $p$. On tame inertia absolute arithmetic Frobenius
acts by $p$-th power. The selected period Frobenius agrees with it
on the unramified constants of $K'$, so their difference lies in
$I G_{K'}$. Conjugation by this difference fixes every character
of that group trivial on $G_{K'}$. The period of each character is at most
$|\Sigma|\le N<p$, hence divides $a$. Its values therefore have order
dividing $e$. The map $\omega(g)=g(\pi_0)/\pi_0$ identifies the
cyclic quotient of tame inertia of order $e$ with $\mu_e$:
$X^e+p$ is Eisenstein over $\mathbb Q_{p^a}$ and splits after adjoining
$\pi_0$. Consequently the restricted characters are uniquely
$\omega^c$ for $c$ in a set $A\subset\mathbb Z/e$. We have $0\in A$,
$pA=A$, and $|A|\le N$. For $R=R(A)$, the largest period under $p$
of $c-d$ with $c,d\in A$, injectivity on restricted characters
identifies $R$ with the original ratio period. For two different
orbits of sizes $i,j$, the period divides $\mathrm{lcm}(i,j)\le
ij\le N^2/4$. Within the same orbit it divides its length, at most
$N-1$ outside the fixed point $0$. Thus $1\le R$ and $4R\le N^2$.

For either representation $V$, pass to $K''=K_tK'$. The crystalline
period module there is obtained from that over $K'$ by extending
unramified constants (the old periods already supply full dimension).
Inertia of $K''/K_t$ fixes those old periods by the character
decomposition and $\omega|_{I\cap G_{K_t}}=1$, and fixes the new
unramified constants. The remaining action of
$\mathrm{Gal}(K''/K_t)$ is semilinear through its residue-field
quotient. Galois descent for vector spaces gives a full-dimensional
space of crystalline periods over $K_t$. For $g\in I$ the operator
$\prod_{c\in A}(g-\omega(g)^c)$ kills the old periods by their
decomposition, hence the new periods by scalar extension and restriction.
The kernel lattice $T_0=\ker g$ is free and stable. Its rational
space inherits crystallinity, the annihilating polynomial and jumps in
$U$ by crystalline subobject exactness and strictness of de Rham
filtrations, as in T84L1. This produces the three good lattices for
the same $A,U,K_t$.
% END Q84T2

% BEGIN Q84T3
\paragraph{Q84T3. Exactness at the middle of a realized integral triple.}
Suppose that compatible Breuil--Kisin realizations of the three
lattices and the two maps are given on $K_t$. Write $F,G$ for their
matrices over $\prod_{\tau\in\mathbb Z/a}S$, where $S=O'[[u]]$,
and $E=u^e+p$. The Frobenius matrices of the modules and their
inverses have entries in $S[1/E]$; Frobenius sends $u$ to $u^p$ and
cycles the components. After evaluation in the integral Witt ring
of the tilt with coefficients $O'$, the matrices of the two maps
are identified with the scalar extensions of the lattice maps.
The short exact sequence of free $O'$-modules splits linearly.
Consequently the two evaluated matrices have respectively a left
and a right inverse, their composite is zero, and the ranks add.
The integral evaluation is injective. One can see this on each
unramified component from the $u$-expansion: modulo the coefficient
uniformizer the map evaluates $u$ at the element
$(\pi_j\bmod p)_j$ of positive finite tilt valuation, while a
nonzero residue constant has valuation zero. The first nonzero
coefficient therefore detects the lowest term. Separation in the
coefficient uniformizer gives injectivity before reduction, with
the unramified idempotents separating the components.
In particular $GF=0$.

For each component and each of the two matrices at least one minor
of the required full rank has a coefficient in $O'^\times$.
Indeed the Cauchy--Binet identity for its evaluated one-sided inverse
expresses $1$ in the ideal of the evaluated full minors. If all
these minors in component $\tau$ were divisible by a uniformizer
$\lambda$ of $O'$, multiplication by the image of the component
idempotent would put that nonzero idempotent in $(\lambda)$.
Idempotence would then put it in every $(\lambda^n)$, contradicting
separation. Rank zero uses the empty minor $1$.

Each chosen primitive minor has only finitely many roots in the open
unit disc by Weierstrass preparation over $O'$. Thus only finitely
many points in that disc, across all components, could have deficient
rank. Away from the roots of $E$, the Frobenius intertwining identity
identifies the rank at $u=\beta$ in one component with that at
$u=\beta^p$ in its neighbor. If $0<|\alpha|<1$ were deficient,
among its $p$-th roots at most one could be a root of $E$ since
$(p,e)=1$. Choose a different root and repeat. The increasing norms
give infinitely many deficient points. At zero, the least
$u$-adic orders $d_\tau$ of the full minors are finite.
Over $L'[[u]]$ the Frobenius matrices are invertible at zero.
Change of basis preserves the ideal of the full minors, so
$d_\tau=p d_{\tau-1}$ cyclically. All $d_\tau$ are zero.

The height-one primes of the regular local ring $S$ are generated by
$\lambda$ or, up to units, irreducible distinguished polynomials,
again by preparation. At every such prime the two maps are therefore
split in their respective directions. The kernel of $F$ vanishes
by generic injectivity. For $x\in\ker G$ the rank equality and
$GF=0$ give a unique rational preimage under $F$. Local splitting
makes it integral at every height-one prime, hence integral over
$S$ since $S$ is a normal domain. This proves $\ker G=\mathrm{im}F$.
The finite module $\mathrm{coker}\,G$ is supported at the closed
point in each two-dimensional component, hence has finite length.
These are exactly the three properties of $\mathrm{IsExactTriple}$.
% END Q84T3

% BEGIN Q84T4
\paragraph{Q84T4. The exponent supplied by the monodromy estimates.}
On the common field of Q84T2 put
\[
 w=2(N-1)R,\quad h_L=(N+1)b,\quad c_1=b(R+1),\quad
 c_2=c_1+w,\quad k_*=c_2+b+h_L+w.
\]
Here $M$ is the integral module in each Breuil--Kisin datum,
$Y=(E^{-b}M)^\Delta$, $Y'=u^eY$, and $M^\Delta$ denotes the tame
invariants. Write $\mathrm{Mon}$ for its linear monodromy obstruction
in the principal-parts module. The following calculation is conditional
on the realized tame models with the indicated heights, jumps and
characters, and on three integral monodromy estimates:
\begin{enumerate}
\item Each $c\in Z_f^1(T_2)$ has a representative $y_2\in Y_2$ for
$p^w c$ with $\mathrm{Mon}(y_2)=0$.
\item Every pair $y'_2\in Y'_2$, $m_2\in M_2^\Delta$ with
$\mathrm{Mon}(y'_2+\varphi m_2)=0$ has a pair
$y'\in Y'_1$, $m\in M_1^\Delta$ with zero obstruction and
$G(y'+\varphi m)=p^{k_*}(y'_2+\varphi m_2)$.
\item For such a zero-obstruction pair, $p^{c_1}(y'+\varphi m)$
represents a class of $Z_f^1(T_1)$ under the integral period
comparison. The representation of classes is functorial for $G$,
is invariant under addition of $(\varphi-1)M^\Delta$, and two
classes represented by the same vector differ by an integral coboundary.
\end{enumerate}
Indeed $E^b-p^b=u^e\sum_{i=0}^{b-1}p^i E^{b-1-i}$. Multiplying
$y_2$ by this identity gives $p^b y_2=y'_2+m_2$, with
$m_2=E^b y_2\in M_2^\Delta$ and
$y'_2=-u^e\sum_i p^iE^{b-1-i}y_2\in Y'_2$.
Thus $y'_2+\varphi m_2=p^b y_2+(\varphi-1)m_2$. The principal-part
obstruction vanishes on the added difference, and the integral class
is multiplied by $p^b$. Apply the second and third estimates and compare
the two representatives in the target. This gives
\[
 p^dZ_f^1(T_2)\subseteq g_*Z_f^1(T_1)+B^1(T_2),\qquad
 d=c_1+k_*+b+w=6(N-1)R+(N+5+2R)b.
\]
Since $N\ge2$ and $R\ge1$, $d\le7(NR+b(N+R))$.
Multiplication by the unused power and Q84T1 give the required
constant $C_T=40$ with the original prime guard, and likewise any
$C_T\ge40$ with its stronger guard.

The formal source derives the three estimates from its
Bhatt--Scholze lattice and Fontaine period hypotheses. In the
middle estimate, its first quantitative step is the bound $p^w$ on
the cokernels of $G$ on $M,M^\Delta,Y'$. Finite length in Q84T3
does not specify this exponent. The subsequent steps bound the
obstruction by $p^{c_1}$, pull its integral principal part into
the kernel at cost $p^w$, and solve the kernel monodromy equation
at additional cost $p^{b+h_L}$. Returning to cocycles further
requires the $p^{c_1}$ bound on the comparison vector $\zeta$.
The fuller human proof of these estimates, including the initial
construction of the bounded representative, remains due. The
independent roadmap B8.42bs1d (revision \texttt{e481c74135c4})
also invokes the quantitative construction in its fourth claim.
% END Q84T4

% BEGIN Q84T5
\paragraph{Q84T5. Interpolating the two degree intervals and cancelling their common shift.}
Write $\mathcal L_i$ for the span of the Frobenius columns of $M_i$
over the integral power-series ring. Set $v=u^e$, $B=w+1$ and
$J=H+b+B$. We have $2J+2<p$: use $b<H$, $4R\le N^2$ and
$40(HN^2+N^3)<p$. Fix a tame isotypic index $c\in\mathbb Z/e$.
For $s\in A$ put $d_s=(c-s)\bmod e$ in $[0,e)$. Let $P_1,P_2$
be the monic products with roots $d_s+je$ for $0\le j<J$ and
$J\le j<2J$, respectively. Then
\[
         p^w=U_1 P_1+U_2 P_2,\qquad U_i\in\mathbb Z_p[X].
\]
To see the actual denominators, fix any root $z=d_s+je$ in the union.
In its own residue progression all nonzero differences from $z$ are
units since $2J<p$ and $p\nmid e$. For each $s'\ne s$ let $r\le R$
be the $p$-orbit length of $s'-s$. Then
$m=(d_s-d_{s'})(p^r-1)/e$ is an integer with $|m|<p^r-1$.
For another index $j'$ the difference $z-(d_{s'}+j'e)$ has
the valuation of $m+(j-j')(p^r-1)$, since both $e$ and $p^r-1$
are $p$-units. This integer is nonzero and has absolute value
$<2Jp^r<p^{r+1}$, giving valuation at most $r$. At most one value
of $j'$ in $[0,2J)$ gives positive valuation. Thus the full
Lagrange denominator at $z$ has valuation at most $(N-1)R\le w$.
Interpolate a polynomial $Q$ equal to $p^w$ on the first interval
and zero on the second. The Lagrange formula gives integral
coefficients, and division by the monic $P_i$ gives $P_2\mid Q$,
$P_1\mid p^w-Q$ over $\mathbb Z_p$. This proves the identity.

Here is the remaining module implication. Write $\mathcal H_i$ for the
high-degree submodule of the extended lattice, generated there
by scalar coefficients divisible by $u^{ep}$ in the power-series
field of coefficients. Use the following precise reductions for
the Euler connection $\nabla$: the extended lattice is
$\mathcal L_i+\mathcal H_i$ and is $\nabla$-stable. On each $c$
part, $P_1(\nabla)$ takes $\mathcal L_i$ into
$v^J\mathcal L_i+\mathcal H_i$, and $P_2(\nabla)$ takes
$v^J\mathcal L_i$ into $v^{2J}\mathcal L_i+\mathcal H_i$.
The submodule
$T=v^JG\mathcal L_1+p^B\mathcal L_2+\mathcal H_2$ is
$\nabla$-stable. These are the integral connection inputs needed here.
The tame inertia projector to $c$ preserves the integral lattices:
average with its inverse character over the cyclic group of order
$e$, a $p$-unit. The projectors sum to $1$, preserve $v^J$ and
commute with $G$. The given morphism is horizontal,
$G\nabla_1=\nabla_2G$, so its action commutes with evaluation of
every polynomial in the two connections.

Apply F84R2 with modulus $p^B$ to obtain the coarse congruence $v^J\mathcal L_2\subseteq
G\mathcal L_1+p^B\mathcal L_2$. For any $x\in\mathcal L_2$, project
$v^J x$ to $y$ in its $c$ part and apply the projected coarse
congruence. Write $y=Gz+p^Bz'$ with each term in its $c$ part.
Then $P_1(\nabla)y\in T$: on the source term use the first degree
bound and $G\mathcal H_1\subseteq\mathcal H_2$, and on the error
use stability of the extended lattice. Also $P_2(\nabla)y\in T$:
apply the second degree bound and use
$v^{2J}\mathcal L_2\subseteq v^JG\mathcal L_1+p^B\mathcal L_2$
by the same coarse congruence. The interpolation identity and
stability of $T$ give $p^w y\in T$. Sum the projectors. We obtain
\[
       v^J p^w x=v^J Gz+p^B z'+h,\qquad h\in\mathcal H_2 .
\]
The equality also puts $h$ in $\mathcal L_2$. In the Frobenius-column
basis its coefficients belong to $O'[[u]]$ and are separately
divisible by $u^{ep}$ over $L'[[u]]$. Removing their initial zero
coefficients therefore gives $h=v^ph_0$ with $h_0\in\mathcal L_2$.
Now $v^J(p^wx-Gz-v^{p-J}h_0)=p^Bz'$. In the same basis multiplication
by $v^J$ shifts the coefficient indices by $eJ$, so every coefficient
of the expression in parentheses belongs to $p^BO'$. Finally
$v^{p-J}h_0=v^J(v^{p-2J}h_0)$ belongs to
$G\mathcal L_1+p^B\mathcal L_2$ by the coarse congruence.
Consequently $p^w\mathcal L_2\subseteq G\mathcal L_1+
p^{w+1}\mathcal L_2$. Finite length of the cokernel on $M$ gives
some $p^k$ annihilating it. Applying Frobenius and taking spans
gives the same coarse annihilator on $\mathcal L$. Iterating the
last inclusion $k$ times therefore eliminates the error altogether,
proving $p^w\mathcal L_2\subseteq G\mathcal L_1$ under the stated
connection reductions and the coarse congruence.
% END Q84T5
